% Бүлэг 1

\chapter{Судалгааны хэсэг} % Бүлгийн нэр
\label{Chapter1} % Энэ бүлэг рүү ишлэл хийх бол \ref{Chapter1} командыг ашигла 

%-------------------------------------------------------------------------------
\section*{Системийн зорилго}

Жижиг, дунд үйлдвэрлэлийн газар (бэйкири)-ын өдөр тутмын үйл ажиллагаанд шаардлагатай барааны нөөцийн тооллого, борлуулалттай уялдсан автомат хасалт, бүтээгдэхүүний ангилал бүрийн хадгалах хугацаа (Shelf Life)-д суурилсан автомат хямдралын жагсаалт үүсгэх процессуудыг хамарсан мэдээллийн системийн зохион байгуулалт, шаардлагыг судлах, дүн шинжилгээ хийхэд чиглэгдэнэ. 

%-------------------------------------------------------------------------------
\section{Сэдэв сонгох}
Тайлангийн сэдэв: «Милая» бэйкириний барааны нөөц бүртгэх систем.

%-------------------------------------------------------------------------------
\section{Системийн зорилт}
\begin{enumerate}
    \item Барааны татан авалт бүртгэх: Ажилтан өглөө бүр салбартаа байгаа бүтээгдэхүүн бүрийг тоолж, тоо ширхэгийг системд бүртгэнэ. 
    \item Эхний үлдэгдэл тогтоох: Систем тухайн өдрийн эхний нөөцийн үлдэгдлийг ангилал, нэр төрлөөр тооцож, өмнөх өдрийн эцсийн үлдэгдэлтэй харьцуулан зөрүүг тэмдэглэнэ. 
    \item Борлуулалт хийгдэх: Өдрийн туршид үйлчлүүлэгчид бараа худалдан авахад кассын POS систем эсвэл eBarimt үүснэ. 
    \item Автомат хасалт: Үүссэн борлуулалтын баримт дээрх борлогдсон бүтээгдэхүүний тоо, нэр төрлийг систем автоматаар уншиж, нийт барааны үлдэгдлээс тухайн тоо хэмжээгээр хасна. 
    \item Бодит цагийн үлдэгдэл харагдах: Хэрэглэгч болон менежер өдрийн турш салбар бүрийн барааны бодит үлдэгдлийг систем дээрээс шууд харах боломжтой. 
    \item Хугацааны хяналт (Shelf Life): Систем бүтээгдэхүүний ангилал тус бүрийн хадгалах хугацаа болон бүртгэгдсэн огноо, цагаас хойш өнгөрсөн хугацааг тасралтгүй тооцоолно. 
    \item Автомат хямдрал зарлах: Бүтээгдэхүүний хадгалах хугацаа дуусахад тохируулсан цонх үлдэх үед (жишээ нь, 72 цагийн хадгалах хугацаатай барааны 67 дахь цагаас эхлэн) систем үнийн дүнгийн 30\%-ийн хямдралыг автоматаар зарлаж, хямдралын жагсаалт руу шилжүүлнэ. 
    \item Салбар, ангилалаар бүлэглэх: Систем хямдралд шилжсэн бүтээгдэхүүнийг салбар, ангилалаар нь ялгаж, ажилтан, менежерт харуулна. 
    \item Тайлан үүсгэх: Систем өдөр тутмын болон үечилсэн (долоо хоног, сар) тайланг салбар, ангилалаар, мөн хямдралд шилжсэн ба хаягдал болсон барааны хувь, хэмжээгээр гаргана.
\end{enumerate}

%-------------------------------------------------------------------------------
\section{Сэдэв сонгох болсон үндэслэл}
Одоогийн байдлаар жижиг, дунд бэйкирийн газруудад барааны нөөцийг дэвтэр, Excel хүснэгт зэрэг гар аргаар бүртгэдэг тохиолдол түгээмэл байдаг. Энэ нь дараах бэрхшээлүүдийг үүсгэдэг:

\begin{itemize}
    \item Цаг хугацааны алдагдал болон хүний хүчин зүйлийн алдаа: Өглөөний тооллогыг гар аргаар бүртгэхэд цаг их шаардагдаж, шивэлтийн алдаа гарах магадлал өндөр байдаг.
    \item Системийн ба бодит үлдэгдэл зөрөх: Борлуулалтын кассын систем болон нөөцийн бүртгэл хоорондоо автоматаар холбогдоогүй тул бодит үлдэгдэл ба системийн үлдэгдэл зөрдөг.
    \item Хадгалах хугацааны субьектив хяналт: Бүтээгдэхүүн бүрийн ангилал (кремтэй бялуу, талх, жигнэмэг)-аас хамаарсан хадгалах хугацааг гараар тооцоолоход хүндрэлтэй тул хямдралд шилжүүлэх цагийг ажилтан өөрийн үзэмжээр, субьектив байдлаар тодорхойлдог.
    \item Нэгдсэн аналитик дутагдах: Салбар бүрийн үлдэгдлийг нэгтгэн харах, ангилал тус бүрээр алдагдал болон борлуулалтын шинжилгээ хийх боломж хязгаарлагдмал.
    \item Төлөвлөлт ба нийлүүлэлтийн хүндрэл: Захиалга, түүхий эд худалдан авалтыг зөв төлөвлөхөд бодит цагийн мэдээлэл дутагддаг.
\end{itemize}

Дээрх бэрхшээлүүдийг арилгаж, Тооллого $\rightarrow$ Борлуулалт (POS) $\rightarrow$ Хугацаа ба Хямдрал гэсэн гурван процессыг автоматаар холбосон, алдаа багатай, цаг хугацаа хэмнэсэн бүртгэлийн системийг судлан боловсруулах шаардлагатай гэж үзэн энэ сэдвийг сонгосон.

%-------------------------------------------------------------------------------
\section{Программын судалгаа}
\subsection{Хэрэглэгчийн талаар мэдээлэл }
Тус системийг ашиглах хэрэглэгчийн бүлгүүд, тэдгээрийн үүрэг хариуцлага болон системд гүйцэтгэх үйлдлүүдийг доорх хүснэгтэд үзүүлэв. 

\begin{table}
    \centering
    \begin{tabular}{p{3cm}p{5cm}p{5.5cm}}
        \toprule
        \textbf{Хэрэглэгчийн төрөл } &
        \textbf{Үүрэг, хариуцлага} &
        \textbf{Системд гүйцэтгэх үйлдэл} \\
        \midrule 
        Салбарын ажилтан /   Худалдагч  &
        Өглөө бүр бүтээгдэхүүнийг тоолж бүртгэх, борлуулалт хийх  &
        Тооллогын мэдээлэл оруулах, бодит үлдэгдэл харах, кассын баримт үүсгэх  \\
        
        Салбарын   менежер  &
        Салбарын нөөц, бүтээгдэхүүний хадгалах хугацаа болон хямдралын мэдээлэлд хяналт тавих &
        Тайлан харах, хугацаа дуусаж буй бараа болон хямдралын жагсаалтыг хянах, баталгаажуулах  \\
        
        Нягтлан бодогч / Санхүү  &
        Кассын баримт, борлуулалтын гүйлгээг санхүүгийн бүртгэлтэй уялдуулах &
        Борлуулалтын болон нөөцийн тайлан татах, зөрүү шалгах, санхүүгийн бүртгэлд тусгах \\
        Ерөнхий менежмент / Захирал  &
        Бүх салбарын нэгдсэн үзүүлэлт, борлуулалтын динамик, алдагдалд шинжилгээ хийх  &
        Нэгдсэн тайлан, аналитик диаграмм харах, шийдвэр гаргах  \\
        
        Системийн администратор  &
        Системийн тасралтгүй ажиллагаа, тохиргоо, хэрэглэгчийн эрх удирдах &
        Хэрэглэгчийн эрх нэмэх/засах, ангилал ба салбар тохируулах \\
        
        Үйлчлүүлэгч  &
        Бүтээгдэхүүн болон хямдралтай барааны мэдээлэл харах  &
        Веб/Апп-аар барааны цэс, хугацаа ба хямдралын мэдээлэл харах, худалдан авалт хийх  \\
        \bottomrule

    \end{tabular}
    \caption{Хэрэглэгчийн талаар мэдээлэл}
    \label{tab:users}
\end{table}

Ажилчид үндсэндээ 5 түвшний хангалтын эрхтэй бөгөөд салбарын ажилтан зөвхөн өөрийн салбарын өгөгдөлтэй ажиллах, харин менежер, санхүүгийн ажилтан болон захирал бүх салбарын нэгтгэсэн мэдээлэлтэй ажиллах боломжтой байхаар зохион байгуулагдана. 

%-------------------------------------------------------------------------------
\section{Системийг судлах}

\subsection{Систем, дэд системүүд}
«Милая» бэйкириний барааны нөөц бүртгэх систем нь дараах 6 үндсэн дэд системээс бүрдэнэ:

\begin{enumerate}
    \item Тооллогын дэд систем: Өглөөний барааны тоо ширхэгийг бүртгэх, эхний үлдэгдэл ба зөрүүг тооцох.
    \item Борлуулалтын дэд систем: Кассын POS систем болон eBarimt-аас борлуулалтын мэдээлэл авах, нөөцийг автоматаар шинэчлэх.
    \item Нөөцийн удирдлагын дэд систем: Нийт үлдэгдэл, салбар/ангилалын хуваарилалт болон бодит цагийн нөөцийг тооцоолох.
    \item Хугацаа, хямдралын удирдлагын дэд систем: Бүтээгдэхүүн бүрийн бүртгэгдсэн цагаас хойших хадгалах хугацааг тооцож, хямдралын цонх эхлэхэд 30\%-ийн хямдралыг автоматаар зарлан жагсаалт руу шилжүүлэх.
    \item Тайлан, дүн шинжилгээний дэд систем: Өдөр, долоо хоног, сарын нэгдсэн тайлан гаргах, график ба аналитик харуулах.
    \item Хэрэглэгч, эрхийн удирдлагын дэд систем: Хэрэглэгчийн нэвтрэлт, хангалтын эрх, салбарын харьяаллыг удирдах.
\end{enumerate}

\subsection{Оролт, боловсруулалт, гаралтын дүрслэл (IPO загвар)}
Дэд систем бүрийн Оролт (Input), Боловсруулалт (Process), Гаралт (Output)-ыг дараах хүснэгтээр харуулав.

\begin{table}
    \centering
    \begin{tabular}{p{2.8cm}p{3.6cm}p{4.2cm}p{3.4cm}}
        \toprule
        \textbf{Дэд систем} &
        \textbf{Оролт (Input)} &
        \textbf{Боловсруулалт (Process)} &
        \textbf{Гаралт (Output)} \\
        \midrule
        Тооллогын дэд систем &
        Ажилтны тоолсон тоо, салбар, огноо, ангилал &
        Эхний үлдэгдлийг өмнөх өдрийн эцсийн үлдэгдэлтэй харьцуулж, зөрүүг тооцох &
        Тухайн өдрийн эхний нөөцийн баталгаажсан бүртгэл \\

        Борлуулалтын дэд систем &
        POS / Кассын баримтын мэдээлэл (Нэр, тоо, үнэ) &
        Борлогдсон тоо хэмжээг нөөцийн сангаас хасах тооцоолол хийх &
        Борлогдсон барааны түүх, шинэчлэгдсэн үлдэгдэл \\

        Нөөцийн удирдлагын дэд систем &
        Эхний үлдэгдэл, Борлуулалтын гүйлгээ &
        Бодит үлдэгдэл = Эхний үлдэгдэл $-$ Борлуулалт формулаар бодит цагт тооцоолох &
        Бодит цагийн нөөцийн үлдэгдэл (салбар, ангилалаар) \\

        Хугацаа, хямдралын удирдлага &
        Бараа бүртгэсэн огноо/цаг, хадгалах хугацаа, хямдралын дүрэм &
        Бүртгэгдсэн цагаас өнгөрсөн хугацааг тооцож, хямдралын цонхонд заасан 30\%-ийн хямдралыг тооцох &
        Автоматаар хямдарсан барааны жагсаалт (салбар, ангилалаар) \\

        Тайлан, дүн шинжилгээ &
        Тооллого, борлуулалт, хямдрал, хаягдлын өгөгдөл &
        Өгөгдлийг хугацаа болон салбараар нэгтгэн аналитик тооцоолол хийх &
        Өдөр/сар бүрийн нэгдсэн тайлан, график ба статистик \\
        \bottomrule
    \end{tabular}
    \caption{IPO загвар}
    \label{tab:ipo}
\end{table}

\subsection{Системд байгаа асуудлууд}
Одоогийн гар бүртгэлд суурилсан нөхцөл байдлаас тодорхойлогдсон гол асуудлууд:

\begin{itemize}
    \item Бүртгэлийн зөрүү: Тооллого ба борлуулалтын мэдээлэл хооронд автомат холбоос байхгүйгээс үүдэн үлдэгдлийн зөрүү гардаг.
    \item Төвлөрсөн хяналтгүй байдал: Салбар хоорондын мэдээлэл нэгтгэгдэхгүй тул удирдлагад бодит цагийн дүр зураг харагдахгүй.
    \item Худалдааны алдагдал: Бараа бүрийн бүртгэгдсэн цагаас хойш өнгөрсөн хугацааг гараар тооцох боломжгүй тул хадгалах хугацааны дүрмийг цаг тухайд нь мөрдөж, хямдрал зарлах боломжгүй байдаг.
    \item Хугацаа алдалт: Тайлан гаргах ажил гараар хийгддэг тул алдаа, хугацаа алдалт их үүсдэг.
    \item Аюулгүй байдал: Мэдээллийн аюулгүй байдал болон хэрэглэгчийн эрхийн зааг ялгаа тодорхой бус.
\end{itemize}

%-------------------------------------------------------------------------------
\section{Хууль, дүрэм журам}
Монгол Улсад онлайн пицца захиалгыг дангаар зохицуулсан тусгай хууль одоогоор байхгүй. Харин цахим худалдаа, хэрэглэгчийн эрх, хувийн мэдээлэл, татвар болон хүнсний аюулгүй байдлын хэд хэдэн хууль хамт үйлчилнэ.


\subsection{Хамаарах хууль, заалт}
\begin{table}
    \centering
    \begin{tabular}{p{4cm}p{4.5cm}p{5.5cm}}
\toprule
\textbf{Хууль, эрх зүйн акт} &
\textbf{Гол заалт} &
\textbf{Системд тавигдах шаардлага} \\
\midrule

\textbf{Хэрэглэгчийн эрхийг хамгаалах тухай хууль} &
4.1.2, 5.1, 7.1--7.2 &
Бүтээгдэхүүний нэр, орц, найрлага, хэмжээ, үнэ, хадгалах болон
хэрэглэх мэдээллийг үнэн зөв харуулна. \\

\textbf{Хүний хувийн мэдээлэл хамгаалах тухай хууль} &
8, 18, 20, 22 дугаар зүйл &
Нэр, утас, хаяг, байршил зэрэг мэдээллийг зөвшөөрөлтэй цуглуулж,
хамгаална. \\

\textbf{Хүнсний бүтээгдэхүүний аюулгүй байдлыг хангах тухай хууль} &
3.1, 9, 10 дугаар зүйл &
Бүтээгдэхүүний орц, түүхий эдийн гарал, хадгалалт, үйлдвэрлэл, хүргэлтийн
бүртгэл хөтөлнө. \\

\textbf{Хүнсний тухай хууль} &
10 дугаар зүйл болон холбогдох заалт &
Хүнсний чиглэлийн үйл ажиллагаа эрхлэгч аюулгүй, шаардлага хангасан
хүнс нийлүүлнэ. \\

\textbf{Нэмэгдсэн өртгийн албан татварын тухай хууль} &
17.3--17.5 &
Борлуулалт бүрийг бүртгэж, төлбөрийн баримт буюу eBarimt олгоно. \\

\textbf{Цахим гарын үсгийн тухай хууль} &
3 дугаар зүйл &
Албан ёсны цахим баримт бичиг, гэрээнд шаардлагатай бол цахим эсвэл
тоон гарын үсэг хэрэглэнэ. \\

\textbf{Хоол үйлдвэрлэл, үйлчилгээний газрын эрүүл ахуй, халдвар
хамгааллын үлгэрчилсэн дүрэм} &
1.3--1.7 болон холбогдох бүлгүүд &
Гал тогоо, тоног төхөөрөмж, ажилтны ариун цэвэр, хадгалалт,
цэвэрлэгээ, дотоод хяналтыг зохицуулна. \\

\bottomrule
        
    \end{tabular}
    \caption{Хуулийн судалгаа}
    \label{tab:laws}
\end{table}

Хэрэглэгч барааны талаар бодит мэдээлэл авах эрхтэй бөгөөд үйлдвэрлэгч нэр, хаяг, бүтээгдэхүүний орц, найрлага, нэр төрөл, үнэ, хэмжээ зэрэг мэдээллийг өгөх үүрэгтэй.

\subsection{Хувийн мэдээллийн шаардлага}
Онлайн систем дараах мэдээллийг цуглуулж болно:
\begin{itemize}
    \item захиалагчийн нэр;
    \item утасны дугаар;
    \item хүргэлтийн хаяг;
    \item байршлын мэдээлэл;
    \item захиалгын түүх;
    \item төлбөрийн төлөв.
\end{itemize}

Мэдээлэл цуглуулахаас өмнө хэрэглэгчид зорилго, цуглуулах мэдээллийн жагсаалт, хадгалах хугацаа, бусдад дамжуулах эсэх болон зөвшөөрлөө цуцлах аргыг мэдээлнэ. Зөвшөөрлийг цаасан эсвэл цахим хэлбэрээр авч болно.

Системд дараах хамгаалалт шаардлагатай:
\begin{itemize}
    \item хэрэглэгчийн зөвшөөрлийг бүртгэх;
    \item нууцлалын бодлого харуулах;
    \item ажилтны эрхийг албан тушаалаар хязгаарлах;
    \item хүргэгчид зөвхөн шаардлагатай нэр, утас, хаягийг харуулах;
    \item нууц үгийг хамгаалалттай хадгалах;
    \item мэдээллийн өөрчлөлтөд аудитын лог үүсгэх;
    \item мэдээлэл алдагдсан үед авах арга хэмжээний төлөвлөгөөтэй байх;
    \item хэрэглэгч мэдээллээ засах, устгуулах хүсэлт гаргах боломжтой байх.
\end{itemize}

\subsection{Хүнсний аюулгүй байдал}

Хууль нь үйлдвэрлэх, хадгалах, худалдах, үйлчлэх бүх үе шатанд үйлчилнэ.

Системд дараах бүртгэл байвал зохино:
\begin{itemize}
    \item орц, түүхий эдийн нийлүүлэгч;
    \item хадгалалтын хугацаа, температур;
    \item Бүтээгдэхүүн бэлтгэсэн цаг;
    \item гомдол, бүтээгдэхүүн буцаалт;
    \item аюултай бүтээгдэхүүнийг татан авсан бүртгэл.
\end{itemize}

Хүнсний үйл ажиллагаа эрхлэгч нь түүхий эдийн ул мөрийг мөрдөн тогтоох, дотоод хяналтын баримт хадгалах, ажилтнаа эрүүл ахуйн сургалт болон эрүүл мэндийн үзлэгт хамруулах, бүтээгдэхүүнийг эрүүл ахуйн шаардлагын дагуу хадгалж тээвэрлэх үүрэгтэй.

\subsection{Төлбөр ба eBarimt}

НӨАТ суутган төлөгч бол систем:
\begin{itemize}
    \item Төлбөр хийсэн тухай бүр борлуулалтыг бүртгэнэ.
    \item Бэлэн мөнгө, карт, QPay зэрэг төлбөрийн төрлийг хадгална.
    \item Хөнгөлөлтийн өмнөх болон дараахах үнийг ялгаж харуулна.
    \item Худалдаа бүрд eBarimt үүсгэнэ.
    \item Баримтын мэдээллийг борлуулалтын нэгдсэн системд хуульд заасан хугацаанд илгээнэ.
    \item Цуцлагдсан болон буцаагдсан төлбөрийг тусад нь бүртгэнэ.
\end{itemize}

НӨАТ-ын тухай хуулийн 17.3-т борлуулалт бүрд төлбөрийн баримт олгох, мэдээллийг гурван хоногийн дотор нэгдсэн системд илгээхээр заасан.


\subsection{Хамаарах стандартууд}

\begin{table}
    \centering
    \begin{tabular}{p{5cm}p{10cm}}
        \toprule
\textbf{Стандарт} &
\textbf{Хэрэглэх хүрээ} \\
\midrule

\textbf{MNS 4946:2019} --- Хоол үйлдвэрлэл, үйлчилгээний газарт
тавих шаардлага &
Пиццаны газар, гал тогоо, тоног төхөөрөмж, үйлчилгээ, ажилтны нөхцөл \\

\textbf{MNS ISO 22000:2019} --- Хүнсний аюулгүй байдлын менежментийн
тогтолцоо &
Хүнсний эрсдэлийн удирдлага, HACCP, ул мөр, хяналтын тогтолцоо \\

\textbf{MNS CAC RCP 1} --- Хүнсний эрүүл ахуйн ерөнхий зарчим &
Түүхий эдээс хэрэглэгч хүртэлх эрүүл ахуйн зохистой дадал \\

\textbf{MNS 6648:2016} --- Хүнсний сав, баглаа боодлын шошгололтод
тавих шаардлага &
Савласан ундаа, соус болон бусад савласан бүтээгдэхүүний шошго \\

\textbf{MNS 7032:2024} --- Түргэн хөргөх, хөлдөөсөн хоол, бэлэн
бүтээгдэхүүний шаардлага &
Хөлдөөсөн пицца эсвэл урьдчилан бэлтгэсэн бүтээгдэхүүн борлуулах үед \\

\bottomrule
    \end{tabular}
    \caption{Хамаарах стандартууд}
    \label{tab:standards}
\end{table}

MNS 4946:2019 нь хоол үйлдвэрлэл, үйлчилгээний газрын үндсэн салбарын стандарт юм. MNS ISO 22000:2019 нь хүнсний аюулгүй байдлын удирдлага болон HACCP тогтолцоонд тавих шаардлагыг тодорхойлдог.

\subsection{Системд заавал тусгах мэдээлэл}

Захиалга баталгаажуулахын өмнө хэрэглэгчид дараахыг харуулна:

\begin{enumerate}
    \item байгууллагын нэр, хаяг, холбоо барих мэдээлэл;
    \item пиццаны нэр, хэмжээ, орц, харшил үүсгэгч;
    \item үндсэн үнэ;
    \item хөнгөлөлтийн хувь ба нөхцөл;
    \item хүргэлтийн төлбөр;
    \item төлөх нийт дүн;
    \item төлбөрийн хэлбэр;
    \item хүргэлтийн хаяг ба тооцоолсон хугацаа;
    \item захиалга цуцлах, буцаалт хийх нөхцөл;
    \item гомдол гаргах суваг;
    \item нууцлалын бодлого.
\end{enumerate}

Анхаарах нь: Үндэсний стандарт бүр автоматаар заавал мөрдөх шинжтэй биш. Харин хууль, техникийн зохицуулалт, тусгай дүрэм, гэрээ эсвэл зөвшөөрлийн нөхцөлд эш татсан бол заавал мөрдөх шаардлага үүсэж болно. Иймээс үйл ажиллагаа явуулах газрын дүүрэг, хүнсний үйлдвэрлэлийн төрөл болон борлуулах бүтээгдэхүүнээр нь нэмэлт шаардлагыг дахин шалгах хэрэгтэй. Энэ жагсаалтыг 2026 оны 9 дүгээр сарын 22-ны байдлаар албан эх сурвалжид тулгуурлан нэгтгэв.

%-------------------------------------------------------------------------------
\section{Нэр томъёоны тайлбар}
	\begin{longtable}{|p{0.5cm}|p{4cm}|p{8.5cm}|}
        \caption{Нэр томъёоны тайлбар}
        \label{table:1} \\
            \hline
                № & Тодорхойлолт & Тайлбар \\ 
            \hline
            \endfirsthead
            \hline
                № & Тодорхойлолт & Тайлбар \\ 
            \hline
            \endhead
                1 & Нөөц (Inventory) & Тухайн цаг мөчид салбарт эсвэл агуулахад байгаа бүтээгдэхүүний тоо хэмжээ. \\ 
            \hline
                2 & Тооллого (Stock count) & Ажилтны биечлэн тоолж, системд баталгаажуулан бүртгэх үйлдэл. \\ 
            \hline
                3 & POS / Кассын систем & Point of Sale буюу борлуулалтын цэг дээр худалдан авалтыг бүртгэх систем. \\ 
            \hline
                4 & Автомат хасалт & Борлуулалт хийгдэхэд систем нөөцийг өөрөө шинэчлэн бууруулах үйлдэл. \\ 
            \hline
                5 & Хадгалах хугацаа (Shelf Life) & Бүтээгдэхүүн системд бүртгэгдсэн цагаас хойш чанараа алдахгүй байх хугацаа. \\ 
            \hline
                6 & Хямдралын цонх & Хадгалах хугацаа дуусахаас өмнөх, 30\%-ийн хямдрал автоматаар зарлагддаг сүүлийн хугацаа. \\ 
            \hline
                7 & Автомат хямдрал (30\%) & Бараа хямдралын цонхонд орсны дараа системээс автоматаар тооцдог үнийн бууралт. \\ 
            \hline
                8 & Ангилал (Category) & Бүтээгдэхүүнийг төрлөөр нь бүлэглэсэн бүлэг (талх, бялуу, жигнэмэг гэх мэт). \\ 
            \hline
                9 & Салбар (Branch) & Бэйкирийн үйлчилгээ явуулдаг тодорхой байршил, дэлгүүр. \\ 
            \hline
                10 & Дэд систем (Subsystem) & Нийт системийг бүрдүүлэгч, тодорхой бие даасан функц гүйцэтгэдэг хэсэг. \\ 
            \hline
    \end{longtable}

%-------------------------------------------------------------------------------
\section{Лавлах материал}
\begin{itemize}
    \item Барааны нөөцийн удирдлагын систем (Inventory Management System): Ерөнхий онолын үндэслэл, зарчим.
    \item Жижиг, дунд үйлдвэрлэлийн газрын нөөцийн бүртгэлийн стандарт: Санхүүгийн тайлагналын журам ба кассын баримттай холбоотой заалтууд.
    \item Системийн шинжилгээ ба зохиомж (System Analysis and Design): Хичээлийн үндсэн сурах бичиг, лекцийн материал.
    \item Мэдээллийн санг зохион байгуулах онол: Оролт-Боловсруулалт-Гаралт (Input--Process--Output) загварчлал.
    \item Жижиг бизнесийн POS систем: Барааны бүртгэлийн зах зээлд байгаа ижил төстэй програм хангамжуудын судалгаа.
\end{itemize}

